\documentclass[12pt]{article}

\usepackage[T1]{fontenc}
\usepackage[utf8]{inputenc}
\usepackage{amsmath}
\usepackage{amssymb}
\usepackage{array}
\usepackage{booktabs}
\usepackage{geometry}
\usepackage{setspace}
\usepackage{threeparttable}
\usepackage{natbib}
\usepackage[hyphens]{url}
\usepackage{xurl}
\usepackage{hyperref}

\geometry{left=28mm,right=30mm,top=28mm,bottom=28mm}
\onehalfspacing
\setlength{\bibsep}{0pt plus 0.3ex}
\setlength{\emergencystretch}{3em}
\providecommand{\sym}[1]{\textsuperscript{#1}}
\renewcommand{\thepage}{OA-\arabic{page}}
\renewcommand{\thesection}{OA-\Alph{section}}
\renewcommand{\thesubsection}{\thesection.\arabic{subsection}}
\renewcommand{\thetable}{OA-\arabic{table}}
\renewcommand{\theequation}{OA-\arabic{equation}}
\renewcommand{\refname}{References for the Online Appendix}
\hypersetup{
    colorlinks=true,
    linkcolor=black,
    citecolor=black,
    filecolor=magenta,
    urlcolor=blue,
    pdftitle={Online Appendix: The Price of Exclusion},
    pdfauthor={Darcio Genicolo-Martins}
}

\InputIfFileExists{../output/values.tex}{}{%
  \PackageError{online_appendix_v8}{Missing generated macro registry}{Run the analysis pipeline before compiling the appendix.}}

\begin{document}

\begin{center}
{\Large \textbf{Online Appendix for}}\\[0.5em]
{\large \textbf{The Price of Exclusion: SME Set-Asides in Public Procurement}}\\[0.5em]
Darcio Genicolo-Martins
\end{center}

\noindent This appendix reports the checks behind the identification and
policy claims in the paper: institutional timing, reduced-form validation,
structural cost recovery, mechanism robustness, coordination screens, and
welfare calculations. The main text keeps only the evidence needed for the
argument. The appendix records the diagnostics that support, qualify, or bound
that interpretation.
Appendix~OA-A documents the institutional setting; OA-B reports reduced-form
diagnostics; OA-C covers structural cost recovery; OA-D reports mechanism and
coordination robustness; and OA-E gives welfare and policy-design details.

\appendix

\section{Institutional Details}

\subsection{Legal timeline}

Brazil's federal SME procurement framework begins with \LCSmeStatute, which
defined the SME categories used in procurement policy. The statute distinguishes
MicroEmpresas, with annual revenue up to R\$\MERevCeilingBRL, from Empresas de
Pequeno Porte, with annual revenue up to R\$\EPPRevCeilingBRL{} after the
\LCRevenueAct{} update. \LCSetAsideAct{} made SME-only procurement mandatory
for eligible purchases below R\$\setAsideThreshBRL, but implementation depended
on how procurement authorities interpreted the threshold.

In S\~ao Paulo, the relevant interpretation changed from a purchase-notice view
to an item-by-item view. Under the restrictive view, a purchase notice whose
total value exceeded the statutory threshold was not treated as SME-only even
when individual items were below the threshold. The item-level interpretation
made many more items eligible for SME-only treatment. BEC communicated the
technical implementation through Comunicados \comunicadoBECone{} and
\comunicadoBECtwo, while operational take-up in Group~65 occurs around
\policyCutoffMonth. I therefore treat \policyCutoffMonth{} as the empirical
cutoff: it is the mass-adoption date, not the first legal moment at which SME
preferences existed.

\begin{table}[!htbp]
\centering
\begin{threeparttable}
\caption{Institutional timeline}
\label{tab:app_timeline}
\footnotesize
\setlength{\tabcolsep}{3pt}
\begin{tabular}{>{\raggedright\arraybackslash}p{0.23\textwidth}c>{\raggedright\arraybackslash}p{0.56\textwidth}}
\toprule
Instrument & Date & Role in the empirical design \\
\midrule
\LCSmeStatute & 2006 & Federal SME statute; defines SME procurement treatment and preference tools. \\
\leiEstadualSp & 2008 & S\~ao Paulo state SME procurement law. \\
\LCSetAsideAct & 2014 & Makes SME-only procurement mandatory for purchases below R\$\setAsideThreshBRL. \\
TCE-SP eTC \eTCRestrictive & 2015 & Illustrates the restrictive interpretation of the threshold. \\
PGE-SP Parecer \pareceNum & 2017 & Supports the item-level interpretation used in BEC implementation. \\
BEC Comunicados \comunicadoBECone--\comunicadoBECtwo & 2017 & Platform communication enabling SME-only and mixed-exclusivity functionality. \\
\policyCutoffMonth & 2018 & Mass take-up of SME-only treatment in Group~65; empirical cutoff. \\
TCE-SP eTC \eTCExpansive & 2018 & Later ruling consistent with the expansive item-level interpretation. \\
\bottomrule
\end{tabular}
\begin{tablenotes}\footnotesize
\item The table separates legal authority from empirical timing. The design
uses the observed mass take-up in Group~65 rather than treating any
single legal document as the treatment date.
\end{tablenotes}
\end{threeparttable}
\end{table}

\subsection{Platform, products, and sample}

The empirical setting is the BEC electronic procurement platform, created under
Decrees \decreeBECOne{} and \decreeBECTwo{} and later used for Preg\~ao
auctions under Decree \decreePregao. The extract covers \nDataMonths{} months,
from \windowStart{} through \windowEnd, and contains \itemTransactionsCount{}
million item-level transactions and \purchaseOffersCount{} purchase offers.
Group~65 accounts for \groupSixtyfivePctTx{} percent of platform transactions
in this window, which is why the item-level reinterpretation has substantial
empirical bite there.

The structural sample restricts attention to Preg\~ao auctions with at least
\filterMinBidders{} bidders and normalized costs below \(c_\epsilon =
\filterCepsUpper\). The upper-cost filter removes only \nFilteredOutPct{}
percent of bids, so the main results are not driven by discarding a large
fraction of the auction data. The main sample contains \nPreAuctions{}
pre-period auctions and \nPostAuctions{} post-period auctions; pharmaceutical
Preg\~ao auctions account for \nPharmaPregaoAuc{} auctions.

\section{Reduced-Form Evidence}

\subsection{Balance and scope of the comparison}

The reduced form is a timing, sign, and scale check for the structural
exercise, not the source of the decomposition. This distinction matters because
Group~65 is an exposed product group rather than a randomized treated unit. In
pre-period balance checks, Group~65 contributes \balanceNGsixfiveItems{} items
and the pooled controls contribute \balanceNCtrlItems{} items across
\balanceNCtrlGroups{} never-treated groups. The largest absolute standardized
difference is \balanceMaxStdDiff; \balanceNVarsAboveTen{} of
\balanceNVarsTotal{} variables exceed 0.10 in absolute value and
\balanceNVarsAboveTwentyFive{} exceed 0.25. These diagnostics justify the main
text's conservative use of the DiD estimates. Table~\ref{tab:app_balance}
reports the covariates behind this summary.

\begin{table}[!htbp]
\centering
\begin{threeparttable}
\caption{Pre-treatment covariate balance: Group~65 versus never-treated controls}
\label{tab:app_balance}
\scriptsize
\setlength{\tabcolsep}{3pt}
\begin{tabular}{lcccccc}
\toprule
 & \textbf{Group 65} & \multicolumn{4}{c}{\textbf{Control-group mean distribution}} & \textbf{Std. diff.} \\
\cmidrule(lr){3-6}
Variable & Mean & P25 & Median & P75 & Mean & Imbens--Rubin \\
\midrule
Log reference price & 4.036 & 4.057 & 5.385 & 6.874 & 5.657 & -0.219 \\
Bids per item & 7.396 & 8.160 & 9.982 & 13.748 & 11.486 & -0.176 \\
Firms per item & 2.784 & 3.203 & 3.681 & 4.347 & 3.765 & -0.599 \\
Share Preg\~ao & 0.645 & 0.358 & 0.471 & 0.805 & 0.559 & 0.445 \\
Share items below R\$80k & 0.966 & 0.886 & 0.964 & 0.992 & 0.907 & -0.037 \\
Share SME bids & 0.000 & 0.021 & 0.075 & 0.104 & 0.068 & -0.567 \\
Share S\~ao Paulo-state bidders & 0.704 & 0.719 & 0.794 & 0.835 & 0.768 & -0.331 \\
Log bidder-buyer distance & 4.126 & 3.579 & 3.858 & 4.197 & 3.853 & 0.160 \\
Final/reference price & 0.683 & 0.595 & 0.626 & 0.661 & 0.641 & 0.010 \\
\bottomrule
\end{tabular}
\begin{tablenotes}\footnotesize
\item Sample: pre-period items in the DiD panel, restricted to auctions with at
least two bidders. Group~65 contributes \balanceNGsixfiveItems{} items; the
\balanceNCtrlGroups{} never-treated control groups contribute
\balanceNCtrlItems{} items. Standardized differences are Imbens--Rubin
normalized differences \citep{imbens2015}. The table is not a parallel-trends
test; it documents cross-group composition that the fixed-effects design must
absorb.
\end{tablenotes}
\end{threeparttable}
\end{table}

The imbalances in Table~\ref{tab:app_balance} are taken seriously.
They do not by themselves test the identifying assumption. The relevant
question is whether Group~65 had differential pre-trends or anticipatory breaks
relative to the comparison groups, which is why the placebo and alternative-DiD
diagnostics below are reported explicitly.

\begin{table}[!htbp]
\centering
\begin{threeparttable}
\caption{Reduced-form diagnostics with reported standard errors}
\label{tab:app_reduced_form}
\scriptsize
\setlength{\tabcolsep}{3pt}
\begin{tabular}{>{\raggedright\arraybackslash}p{0.36\textwidth}cc>{\raggedright\arraybackslash}p{0.31\textwidth}}
\toprule
Diagnostic & Estimate & Standard error & Sample / notes \\
\midrule
TWFE, headline estimator & \didAltTwfeEst{} & \didAltTwfeSe{} & Item-clustered; \structuralWindowM-month window \\
BJS imputation estimator & \didAltBjsEst{} & \didAltBjsSe{} & Robust imputation; post-treatment ATT \\
Callaway--Sant'Anna estimator & \didAltCsEst{} & \didAltCsSe{} & Never-treated controls; simple post ATT \\
Phased-adoption check, full sample & \phasedDidFull{} & \phasedDidSe{} & Platform enablement period retained \\
Phased-adoption check, excluding enablement months & \phasedDidTrunc{} & \phasedDidSe{} & Months \phasedExcludedFrom--\phasedExcludedTo{} excluded \\
\bottomrule
\end{tabular}
\begin{tablenotes}\footnotesize
\item Estimates use the paper's sign convention. Negative coefficients mean
the treated Group~65 price/reference discount narrows after SME-only adoption,
so prices rise relative to the pre-policy benchmark. The headline TWFE row uses
item fixed effects and item-clustered standard errors. The modern DiD rows are
reported because they discipline the reduced-form validation rather than because
the paper relies on their magnitude for the structural decomposition.
\end{tablenotes}
\end{threeparttable}
\end{table}

\begin{table}[!htbp]
\centering
\begin{threeparttable}
\caption{Placebo timing checks on log prices}
\label{tab:app_placebo}
\footnotesize
\setlength{\tabcolsep}{6pt}
\begin{tabular}{lccc}
\toprule
Fake cutoff & Estimate & Standard error & Observations \\
\midrule
September 2017 & \placeboSeventeenSep{} & 0.0182 & 311,883 \\
March 2017 & \placeboSeventeenMar{} & 0.0132 & 215,611 \\
June 2017 & \placeboSeventeenJun{} & 0.0128 & 237,625 \\
Real-cutoff benchmark range & \placeboLowerNp{} to \placeboHigherNp{} & -- & -- \\
\bottomrule
\end{tabular}
\begin{tablenotes}\footnotesize
\item Placebo regressions use pre-treatment data only, item and month fixed
effects, sealed-bid and log-quantity controls, and item-level clustering. The
placebo estimates are not all zero, so they limit the force of a
reduced-form-only design. Their magnitudes are smaller than the real-cutoff
benchmark, which is why the main text uses the DiD as validation for timing and
sign rather than as the source of the structural price effect.
\end{tablenotes}
\end{threeparttable}
\end{table}

\subsection{Placebo timing and alternative estimators}

Placebo cutoffs before the reform are smaller in magnitude than the real-cutoff
estimates: \placeboSeventeenSep, \placeboSeventeenMar, and
\placeboSeventeenJun. They are not zero in every case, so they do not support a
strong reduced-form-only interpretation. They do support the narrower claim
used in the main text: the largest movement occurs around the institutional
take-up date and has the sign predicted by the set-aside.

Modern DiD estimators reduce precision and attenuate the point estimate. The
BJS imputation estimate is \didAltBjsEst{} and the Callaway--Sant'Anna estimate
is \didAltCsEst, compared with a panel-aggregated TWFE estimate of
\didAltTwfeEst. The reduced form is therefore not asked to carry the structural
price magnitude or a reduced-form welfare calculation. The decomposition uses
the auction model and treats the DiD as external validation for timing and
sign.

\section{Structural Cost Recovery and Estimation}

\subsection{Drop-out bids and winner censoring}

In a Preg\~ao reverse auction, losing firms exit as the price clock falls. Under
the maintained IPV clock interpretation, those exits reveal willingness to
supply at the auction scale. The winner is censored because the winning firm
does not need to reveal the lowest price at which it would still supply. The
baseline estimator treats the winning price as a tight upper-bound observation;
the appendix checks a losers-only estimator and a Turnbull NPMLE
\citep{turnbull1976} that treats the winning observation as left-censored.

\begin{table}[!htbp]
\centering
\begin{threeparttable}
\caption{Winner-censoring and structural cost-recovery checks}
\label{tab:app_cost_recovery}
\begin{tabular}{lcccc}
\toprule
Estimator & \multicolumn{2}{c}{Net price effect} & \multicolumn{2}{c}{Deviation from baseline} \\
\cmidrule(lr){2-3}\cmidrule(lr){4-5}
 & Non-pharma & Pharma & Non-pharma & Pharma \\
\midrule
All bidders, baseline & \costRegimeAllNp{} & \costRegimeAllPh{} & -- & -- \\
Losers only & \costRegimeLosersNp{} & \costRegimeLosersPh{} & \costRegimeDevLoNp\% & \costRegimeDevLoPh\% \\
Turnbull NPMLE & \costRegimeTurnNp{} & \costRegimeTurnPh{} & \costRegimeDevTurnNp\% & \costRegimeDevTurnPh\% \\
\bottomrule
\end{tabular}
\begin{tablenotes}\footnotesize
\item Net price effects are \(p_{S_3}-p_{S_1}\), normalized by the buyer
reference price. The largest Turnbull deviation from the baseline is
\turnbullDeviationPct{} percent; in non-pharmaceuticals it is
\turnbullDeviationNpPct{} percent.
\end{tablenotes}
\end{threeparttable}
\end{table}

The rejections in Table~\ref{tab:app_ks} are not treated as failed robustness
checks. They are the empirical reason for using the observed post-policy SME
pool in the baseline and reporting strict primitive invariance as a bounding
exercise. The policy ranking is therefore emphasized where it survives that
choice, especially in non-pharmaceutical markets.

The conclusion is not that censoring is irrelevant. It is that the central
price-formation ranking does not hinge on the baseline treatment of the winner.
The Turnbull estimator still gives exclusion shares of \turnbullShareNp{}
percent in non-pharmaceuticals and \turnbullSharePh{} percent in
pharmaceuticals.

\begin{table}[!htbp]
\centering
\begin{threeparttable}
\caption{Willingness-to-supply quantiles under alternative winner treatments}
\label{tab:app_turnbull_quantiles}
\scriptsize
\setlength{\tabcolsep}{3pt}
\begin{tabular}{lllrrrrrr}
\toprule
 & & & \multicolumn{3}{c}{Median cost \(c_{0.5}\)} & \multicolumn{3}{c}{Upper-quartile cost \(c_{0.75}\)} \\
\cmidrule(lr){4-6}\cmidrule(lr){7-9}
Class & Type & Period & Losers & All & Turnbull & Losers & All & Turnbull \\
\midrule
Non-pharma & SME & Post & 0.922 & 0.861 & 0.840 & 1.215 & 1.097 & 1.117 \\
Non-pharma & SME & Pre & 0.947 & 0.881 & 0.862 & 1.314 & 1.193 & 1.212 \\
Non-pharma & Non-SME & Post & 0.802 & 0.762 & 0.726 & 1.068 & 0.964 & 0.971 \\
Non-pharma & Non-SME & Pre & 0.768 & 0.719 & 0.684 & 1.028 & 0.918 & 0.932 \\
Pharma & SME & Post & 0.860 & 0.799 & 0.781 & 1.078 & 0.966 & 1.002 \\
Pharma & SME & Pre & 0.863 & 0.800 & 0.710 & 1.441 & 1.188 & 1.245 \\
Pharma & Non-SME & Post & 0.674 & 0.657 & 0.633 & 0.814 & 0.774 & 0.765 \\
Pharma & Non-SME & Pre & 0.619 & 0.604 & 0.574 & 0.762 & 0.725 & 0.708 \\
\bottomrule
\end{tabular}
\begin{tablenotes}\footnotesize
\item Losers-only uses only drop-out bids and is upward-biased because it omits
the lowest-cost bidder in each auction. All-bidders treats the winner's final
price as a point observation and is downward-biased if the final price exceeds
true cost. Turnbull NPMLE treats the winner as left-censored at the final price.
\end{tablenotes}
\end{threeparttable}
\end{table}

\subsection{Auction-level heterogeneity and cross-modality comparison}

The baseline removes auction-level scale heterogeneity before simulating
counterfactual bidder pools. The Preg\~ao intraclass correlations range from
\iccNpSme{} for non-pharmaceutical SMEs to \iccPharmaNonSme{} for
pharmaceutical non-SMEs. These magnitudes are comparable to settings where
auction-level unobserved heterogeneity is important: \citet{krasnokutskaya2011}
reports an ICC benchmark of \krasICC{} in highway procurement.

The cross-modality check compares Preg\~ao drop-out recovery with a GPV
first-price recovery from Convite auctions \citep{gpv2000}. The operational
Kolmogorov--Smirnov distance threshold is \ksDistThresh. The cleanest
falsification cell, pharmaceutical non-SMEs in Convite, has distance
\ksConvitePharmaNS. Non-pharmaceutical Preg\~ao is close to the threshold
(\ksNpPregao), while pharmaceutical Preg\~ao before the heterogeneity correction
is farther away (\ksPharmaPregao). This pattern supports the maintained
correction but also motivates the main text's cautious treatment of IPV.

\begin{table}[!htbp]
\centering
\begin{threeparttable}
\caption{Primitive-stability and cross-modality diagnostics}
\label{tab:app_ks}
\scriptsize
\setlength{\tabcolsep}{3pt}
\begin{tabular}{>{\raggedright\arraybackslash}p{0.33\textwidth}cc>{\raggedright\arraybackslash}p{0.36\textwidth}}
\toprule
Diagnostic cell & KS distance & Threshold & Interpretation \\
\midrule
Non-pharma Preg\~ao, raw drop-outs & \ksNpPregao{} & \ksDistThresh{} & Close to the operational threshold. \\
Pharma Preg\~ao, raw drop-outs & \ksPharmaPregao{} & \ksDistThresh{} & Fails the threshold before common-scale cleaning. \\
Non-pharma Preg\~ao, UH-clean & \ksNpPregaoUH{} & \ksDistThresh{} & Moves above the threshold after removing common scale. \\
Pharma Preg\~ao, UH-clean & \ksPharmaPregaoUH{} & \ksDistThresh{} & Remains far from strict primitive invariance. \\
Non-pharma Convite GPV & \ksNpConvite{} & \ksDistThresh{} & Borderline cross-modality comparison. \\
Pharma Convite GPV, non-SME & \ksConvitePharmaNS{} & \ksDistThresh{} & Cleanest validation cell; passes the threshold. \\
Convite pharma non-SME, post-UH & \ksUHinvCheckLo{} & \ksDistThresh{} & Continues to pass after the heterogeneity correction. \\
\bottomrule
\end{tabular}
\begin{tablenotes}\footnotesize
\item The table reports Kolmogorov--Smirnov distances for diagnostics used to
discipline the cost-recovery step. The goal is not to prove primitive
invariance in every cell; the pharmaceutical Preg\~ao cells visibly reject that
strong restriction. The point is to locate where the recovery is most credible
and where the paper must treat policy rankings as conditional.
\end{tablenotes}
\end{threeparttable}
\end{table}

\section{Robustness to the Core Mechanism}

\subsection{Strict invariance and protected-pool composition}

The decomposition separates the effect of lost non-SME competition from the
protected-pool offset. The weakest interpretive link is the offset, because
\(S_3-S_2\) combines additional SME participation with changes in the active
SME cost distribution. The strict-invariance benchmark therefore imposes the
pre-policy SME cost distribution on the post-policy SME pool.

\begin{table}[!htbp]
\centering
\begin{threeparttable}
\caption{Mechanism robustness}
\label{tab:app_mechanism}
\begin{tabular}{lcccc}
\toprule
Specification & \multicolumn{2}{c}{Non-pharma (main result)} & \multicolumn{2}{c}{Pharma (boundary case)} \\
\cmidrule(lr){2-3}\cmidrule(lr){4-5}
 & Net effect & Excl. share & Net effect & Excl. share \\
\midrule
Baseline & \bneEffectTotalNp{} & \bneShareIntensNp\% & \bneEffectTotalPh{} & \bneShareIntensPh\% \\
Strict invariance & \siDeltaTotalNp{} & \strictShareNp\% & \siDeltaTotalPh{} & \strictSharePh\% \\
Turnbull cost recovery & \costRegimeTurnNp{} & \turnbullShareNp\% & \costRegimeTurnPh{} & \turnbullSharePh\% \\
18-month window & \windowTotalEighteenNp{} & \windowShareEighteenNp\% & \windowTotalEighteenPh{} & \windowShareEighteenPh\% \\
12-month window & \windowTotalTwelveNp{} & \windowShareTwelveNp\% & \windowTotalTwelvePh{} & \windowShareTwelvePh\% \\
6-month window & \windowTotalSixNp{} & \windowShareSixNp\% & \windowTotalSixPh{} & \windowShareSixPh\% \\
\bottomrule
\end{tabular}
\begin{tablenotes}\footnotesize
\item Net effects are normalized by the reference price. Exclusion share is
the absolute share of the decomposition attributed to removing non-SMEs from
the price-forming pool. Rows combine different diagnostic families: strict
invariance disciplines the protected-pool composition margin, Turnbull
disciplines winner censoring, and window rows discipline timing sensitivity.
The pharmaceutical boundary case is the scope condition identified in
main-text~\S4.4; welfare-ranking inversion under strict invariance is reported
in Table~\ref{tab:app_welfare_ranking}.
\end{tablenotes}
\end{threeparttable}
\end{table}

The non-pharmaceutical result is stable: the exclusion component remains the
dominant force under baseline, strict invariance, Turnbull censoring, and
alternative windows. The pharmaceutical result is informative but more
conditional. Post-policy SME turnover is larger in pharmaceuticals:
\turnoverPhPctNewFirms{} percent of post-policy SME firms are new relative to
the pre-period pool and these firms account for \turnoverPhPctNewBids{} percent
of post-policy SME bids. In non-pharmaceuticals, new firms account for
\turnoverNpPctNewBids{} percent of post-policy SME bids. Pharmaceuticals are
therefore a boundary case for policy ranking rather than the headline market.

\subsection{Entry counts, filters, and dependence}

The baseline summarizes entry using class-by-period arrival rates. Window
checks show little drift in the non-pharmaceutical exclusion share:
\windowShareDriftNp{} percentage points across the 18-, 12-, and 6-month
windows. The pharmaceutical drift is larger, at \windowShareDriftPh{} points,
again matching the interpretation that pharma is composition-sensitive.

Cost-filter checks also preserve the qualitative result. Tightening the filter
reduces the non-pharmaceutical net effect by \filterTightDropNp{} percent and
the pharmaceutical net effect by \filterTightDropPh{} percent; loosening it
raises the effects by at most \filterLooseUpNp{} percent and
\filterLooseUpPh{} percent, respectively. Allowing within-auction cost
correlation up to \(\rho_c=\ipvRhoMax\) moves the exclusion share by at most
\ipvShareDriftPp{} percentage points and the total effect by at most
\ipvTotalDriftPct{} percent.

\begin{table}[!htbp]
\centering
\begin{threeparttable}
\caption{Filter, window, and dependence sensitivities}
\label{tab:app_filter_dependence}
\scriptsize
\setlength{\tabcolsep}{3pt}
\begin{tabular}{>{\raggedright\arraybackslash}p{0.31\textwidth}cc>{\raggedright\arraybackslash}p{0.39\textwidth}}
\toprule
Sensitivity & Non-pharma & Pharma & Interpretation \\
\midrule
Tight cost filter: net effect & \filterTightNp{} & \filterTightPh{} & Smaller net effect after removing high normalized costs. \\
Baseline filter: net effect & \filterBaseNp{} & \filterBasePh{} & Benchmark filter used in the robustness registry. \\
Tight-filter drop & \filterTightDropNp\% & \filterTightDropPh\% & Attenuation relative to the benchmark filter. \\
Loose-filter increase & \filterLooseUpNp\% & \filterLooseUpPh\% & Maximum increase under the looser filter. \\
Exclusion-share drift across windows & \windowShareDriftNp{} pp & \windowShareDriftPh{} pp & Difference across 18-, 12-, and 6-month windows. \\
Dependence grid & \multicolumn{2}{c}{\(\rho_c \in \ipvRhoGrid\)} & Total-effect drift at most \ipvTotalDriftPct\%; exclusion-share drift at most \ipvShareDriftPp{} pp. \\
\bottomrule
\end{tabular}
\begin{tablenotes}\footnotesize
\item Net effects are \(p_{S_3}-p_{S_1}\), normalized by the reference price.
The dependence sensitivity allows within-auction cost correlation up to
\(\rho_c=\ipvRhoMax\). The table does not require every specification to have
the same level. It asks whether the exclusion-dominant decomposition survives
the main filter, window, and dependence perturbations.
\end{tablenotes}
\end{threeparttable}
\end{table}

\subsection{Coordination screens}

Coordination is a first-order threat because repeated bidder pairs could make
drop-out prices reflect conduct rather than costs. The screens detect
clustering, but the important question for this paper is whether clustering
intensifies after non-SMEs are removed. The Conley-style screen is stable in
non-pharmaceuticals: realized shares are \collConleyRealPreNp{} percent before
and \collConleyRealPostNp{} percent after the cutoff, compared with null means
of \collConleyNullPreNp{} and \collConleyNullPostNp. In pharmaceuticals, the
realized share falls from \collConleyRealPrePh{} to \collConleyRealPostPh{}
percent. Bajari--Ye-style ratios show the same directional pattern: the
non-pharmaceutical ratio falls from \collBYRatioPreNp{} to
\collBYRatioPostNp, and the pharmaceutical ratio falls from
\collBYRatioPrePh{} to \collBYRatioPostPh.

These screens do not prove competitive conduct. They answer the narrower threat
to the decomposition: the main post-policy price increase is not obviously
driven by a new coordination shock among the surviving SME pool.

\begin{table}[!htbp]
\centering
\begin{threeparttable}
\caption{Bidder-pair coordination screens}
\label{tab:app_coordination}
\footnotesize
\setlength{\tabcolsep}{4pt}
\begin{tabular}{lccc}
\toprule
Stratum & Conley realized & Conley null mean & Bajari--Ye \(T_1\) obs/null \\
\midrule
Non-pharma Pre & \collConleyRealPreNp\% & \collConleyNullPreNp\% & \collBYRatioPreNp{} \\
Non-pharma Post & \collConleyRealPostNp\% & \collConleyNullPostNp\% & \collBYRatioPostNp{} \\
Pharma Pre & \collConleyRealPrePh\% & \collConleyNullPrePh\% & \collBYRatioPrePh{} \\
Pharma Post & \collConleyRealPostPh\% & \collConleyNullPostPh\% & \collBYRatioPostPh{} \\
\bottomrule
\end{tabular}
\begin{tablenotes}\footnotesize
\item The Conley-style statistic is the share of auctions in which the closest
pair of UH-clean residual bids lies within a narrow reference-price band. The
null resamples bid sequences within cell while preserving bidder counts. The
Bajari--Ye statistic is the ratio of the observed maximum pair co-bidding count
to its permutation-null mean among repeated bidder pairs. Both screens detect
baseline clustering. The identifying question here is the pre/post direction:
neither statistic rises after the set-aside expands.
\end{tablenotes}
\end{threeparttable}
\end{table}

\subsection{Reference-price validation}\label{oa:pref_validation}

All structural prices are normalized by buyer reference prices, so mechanical
changes in reference-price setting at the cutoff would be a problem. The
direct evidence below shows that $p^{\mathrm{ref}}$ moves only weakly and
that the structural decomposition is invariant to that movement.

The TWFE DiD on $\log p^{\mathrm{ref}}$ under the main specification
(Table~\ref{tab:app_pref_endog}) is \pRefDidEstNp{} (SE \pRefDidSeNp{}) in
non-pharmaceuticals and \pRefDidEstPh{} (SE \pRefDidSePh{}) in
pharmaceuticals---an order of magnitude smaller than the corresponding DiD
on $\log p^{\mathrm{final}}$ in each subset (\pRefSanityFinalEstNp{} and
\pRefSanityFinalEstPh{}). BJS imputation gives slightly larger magnitudes
(\pRefDidBjsEstNp{} NP, \pRefDidBjsEstPh{} PH). Callaway-Sant'Anna estimates
at group-level aggregation have wide standard errors (above 0.3 in each
subset) and are statistically indistinguishable from zero.

\begin{table}[!htbp]
\centering
\begin{threeparttable}
\caption{Reference-price endogeneity: DiD on $\log p^{\mathrm{ref}}$ across estimators and classes}
\label{tab:app_pref_endog}
\footnotesize
\setlength{\tabcolsep}{4pt}
\begin{tabular}{lcccccc}
\toprule
 & \multicolumn{3}{c}{Non-pharma} & \multicolumn{3}{c}{Pharma} \\
\cmidrule(lr){2-4}\cmidrule(lr){5-7}
Estimator & $\hat\delta$ & SE & N & $\hat\delta$ & SE & N \\
\midrule
TWFE                & \pRefDidEstNp{}    & \pRefDidSeNp{}    & \pRefDidObsNp{} & \pRefDidEstPh{}    & \pRefDidSePh{}    & \pRefDidObsPh{} \\
Callaway--Sant'Anna & \pRefDidCsEstNp{}  & \pRefDidCsSeNp{}  & ---             & \pRefDidCsEstPh{}  & \pRefDidCsSePh{}  & --- \\
BJS imputation      & \pRefDidBjsEstNp{} & \pRefDidBjsSeNp{} & ---             & \pRefDidBjsEstPh{} & \pRefDidBjsSePh{} & --- \\
\addlinespace
\multicolumn{7}{l}{\textit{Sanity: same spec on $\log p^{\mathrm{final}}$ (TWFE)}} \\
$\quad$ Coefficient & \pRefSanityFinalEstNp{} & ---              & ---             & \pRefSanityFinalEstPh{} & ---              & --- \\
\addlinespace
\multicolumn{7}{l}{\textit{Heterogeneity by $|\Delta\log p^{\mathrm{ref}}|$ quartile (DV: $\log p^{\mathrm{final}}$)}} \\
$\quad$ Q1 (low)    & \pRefHetQOneNp{}   & \pRefHetQOneSeNp{}   & ---           & \pRefHetQOnePh{}   & \pRefHetQOneSePh{}   & --- \\
$\quad$ Q4 (high)   & \pRefHetQFourNp{}  & \pRefHetQFourSeNp{}  & ---           & \pRefHetQFourPh{}  & \pRefHetQFourSePh{}  & --- \\
\bottomrule
\end{tabular}
\begin{tablenotes}\footnotesize
\item TWFE uses item and month fixed effects with item-clustered standard
errors, mirroring the main DiD specification. Callaway-Sant'Anna aggregates
to group level for tractability; the standard errors are wide. BJS uses
item-level imputation. The sanity row reports the main DiD on
$\log p^{\mathrm{final}}$ in the same NP/PH sample. The quartile rows split
items by absolute change in $\log p^{\mathrm{ref}}$ pre-to-post and re-run
the DiD on $\log p^{\mathrm{final}}$ within each quartile.
\end{tablenotes}
\end{threeparttable}
\end{table}

The internal placebo by quartile of $|\Delta\log p^{\mathrm{ref}}|$ shows
that items where the reference price was stable across the cutoff still
display a price effect of \pRefHetQOneNp{} (NP) and \pRefHetQOnePh{} (PH),
about two-thirds of the class-specific DiD coefficient on
$\log p^{\mathrm{final}}$. A density test at the R\$\,\setAsideThreshBRL{}
item-eligibility threshold fails to detect strategic manipulation
(post-period T = \bunchingDensityTPost{}, p = \bunchingDensityPPost{};
pre-period T = \bunchingDensityTPre{}, p = \bunchingDensityPPre{}), ruling
out the natural buyer-side channel.

Independent corroboration comes from the second-lowest bid in the iterative
Preg\~ao phase. The DiD on $\log b^{(2)}$ delivers \bidSecondDidEstNp{}
(SE \bidSecondDidSeNp{}) in non-pharmaceuticals and \bidSecondDidEstPh{}
(SE \bidSecondDidSePh{}) in pharmaceuticals---constructed from raw bid
units, with no $p^{\mathrm{ref}}$ normalization---and corroborates the
order-statistic mechanism directly. The structural net effects of
\structRatioNp{} percent and \structRatioPh{} percent of the open-regime
price exceed the reduced-form benchmark of \didImpliedRatio{} percent
because the structural exercise simulates the order statistic under fixed
bidder pools while the reduced form averages over take-up, controls, and
product-time variation. That gap is a reason to interpret levels carefully,
not a reason to discard the mechanism.

\section{Welfare Details}

\subsection{Lambda grid}

The main text reports welfare losses at \(\lambda=\lambdaMain\). The welfare
loss is
\[
    L(\lambda)=DWL^{alloc}+\lambda\Delta G,
\]
where \(DWL^{alloc}\) is the allocative loss and \(\Delta G\) is the extra
government outlay. The grid ranges from \welfLambdaRangeLo{} to
\welfLambdaRangeHi{} in increments of \lambdaInc. Moving \(\lambda\) by 0.10
changes the welfare loss by about \welfLossDeltaNpPp{} percentage points in
non-pharmaceuticals and \welfLossDeltaPhPp{} points in pharmaceuticals.

\begin{table}[!htbp]
\centering
\begin{threeparttable}
\caption{Welfare loss across marginal-cost-of-public-funds values}
\label{tab:app_lambda_grid}
\begin{tabular}{lccccc}
\toprule
Class & \(\lambda=0.15\) & \(\lambda=0.20\) & \(\lambda=0.30\) & \(\lambda=0.40\) & \(\lambda=0.45\) \\
\midrule
Non-pharma & \welfLossPctLfifteenNp\% & \welfLossPctLtwentyNp\% & \welfLossPctLthirtyNp\% & \welfLossPctLfortyNp\% & \welfLossPctLfortyfiveNp\% \\
Pharma & \welfLossPctLfifteenPh\% & \welfLossPctLtwentyPh\% & \welfLossPctLthirtyPh\% & \welfLossPctLfortyPh\% & \welfLossPctLfortyfivePh\% \\
\bottomrule
\end{tabular}
\begin{tablenotes}\footnotesize
\item Losses are reported as percentages of the open-regime simulated price.
\end{tablenotes}
\end{threeparttable}
\end{table}

\subsection{Preference grid and welfare weights}

The 10 percent SME price preference is a benchmark instrument. It keeps
non-SMEs inside the auction and tilts allocation through the scoring rule. In
the simulated preference grid from 0 to 30 percent, the 10 percent preference
is the cheapest welfare-dominant preference threshold in the current grid. At
10 percent, the price shift is \vThreeDeltaNp{} in non-pharmaceuticals and
\vThreeDeltaPh{} in pharmaceuticals, while SME win rates rise by
\prefSmeWinGainTenNp{} and \prefSmeWinGainTenPh{} percentage points.

The preference is not a costless substitute for the set-aside. At 30 percent,
the welfare loss remains \welfPrefThirtyNp{} percent in non-pharmaceuticals and
\welfPrefThirtyPh{} percent in pharmaceuticals, but the instrument still
delivers less redistribution than full exclusion. The relevant comparison is
therefore a frontier: how much additional SME surplus, access, or dynamic
capacity does the planner require to justify removing non-SMEs from the
price-forming pool? In the main specification, the implicit SME welfare weights
required to prefer the set-aside to the 10 percent preference are
\welfWeightMainNp{} in non-pharmaceuticals and \welfWeightMainPh{} in
pharmaceuticals. Under strict invariance, the non-pharmaceutical comparison
continues to require a weight above one, while the pharmaceutical comparison
falls to \welfWeightSiPh.

\begin{table}[!htbp]
\centering
\begin{threeparttable}
\caption{Preference grid: price effects, SME win rates, and welfare loss}
\label{tab:app_preference_grid}
\scriptsize
\setlength{\tabcolsep}{3pt}
\begin{tabular}{llrrrrr}
\toprule
Class & Preference & \(\bar p\) & \(\Delta\) vs. \(S_1\) & \% of V0 \(\Delta\) & SME win rate & Welfare loss \\
\midrule
Non-pharma & 0\% & 0.772 & +0.0046 & 2.0 & 17.1\% & 0.14\% \\
Non-pharma & 5\% & 0.768 & +0.0015 & 0.7 & 20.0\% & 0.00\% \\
Non-pharma & 10\% & 0.764 & -0.0028 & -1.2 & 21.4\% & 0.25\% \\
Non-pharma & 15\% & 0.766 & -0.0013 & -0.6 & 24.6\% & 0.00\% \\
Non-pharma & 20\% & 0.775 & +0.0081 & 3.5 & 27.8\% & 1.40\% \\
Non-pharma & 25\% & 0.790 & +0.0227 & 9.7 & 30.2\% & 2.43\% \\
Non-pharma & 30\% & 0.809 & +0.0420 & 17.9 & 34.0\% & 4.99\% \\
Non-pharma & Full set-aside & 1.001 & +0.2342 & 100.0 & 100.0\% & 27.13\% \\
Pharma & 0\% & 0.648 & -0.0061 & -2.1 & 16.2\% & 0.00\% \\
Pharma & 5\% & 0.670 & +0.0161 & 5.5 & 16.5\% & 1.24\% \\
Pharma & 10\% & 0.648 & -0.0061 & -2.1 & 17.6\% & 0.00\% \\
Pharma & 15\% & 0.655 & +0.0009 & 0.3 & 19.6\% & 0.00\% \\
Pharma & 20\% & 0.655 & +0.0015 & 0.5 & 19.8\% & 0.00\% \\
Pharma & 25\% & 0.656 & +0.0019 & 0.7 & 22.1\% & 0.00\% \\
Pharma & 30\% & 0.676 & +0.0219 & 7.5 & 25.7\% & 1.92\% \\
Pharma & Full set-aside & 0.948 & +0.2941 & 100.0 & 100.0\% & 41.99\% \\
\bottomrule
\end{tabular}
\begin{tablenotes}\footnotesize
\item Bayes--Nash simulation with \prefGridB{} Monte Carlo draws per cell.
Preference rate \(k\) scores SME bids by a factor \(1-k\) for winner selection
while paying the actual winning bid. \(\Delta\) is the simulated price effect
relative to \(S_1\), in reference-price units. Welfare loss is
\((DWL^{alloc}+\lambda \Delta G)/p_{S_1}\) at \(\lambda=\lambdaMain\). Entries
that are negative before rounding are reported as 0.00\%, since they lie inside
Monte Carlo noise around a near-zero welfare effect.
\end{tablenotes}
\end{threeparttable}
\end{table}

\begin{table}[!htbp]
\centering
\begin{threeparttable}
\caption{Set-aside versus preference ranking across \(\lambda\)}
\label{tab:app_welfare_ranking}
\footnotesize
\setlength{\tabcolsep}{4pt}
\begin{tabular}{lccccc}
\toprule
 & \multicolumn{2}{c}{Non-pharma} & \multicolumn{3}{c}{Pharma} \\
\cmidrule(lr){2-3}\cmidrule(lr){4-6}
\(\lambda\) & Loss(V0) & Ranking & Loss(V0) & Ranking, main & Ranking, strict inv. \\
\midrule
0.15 & 24.0\% & V3 \(>\) V0 & 38.4\% & V3 \(>\) V0 & V0 \(>\) V3 \\
0.20 & 25.6\% & V3 \(>\) V0 & 40.8\% & V3 \(>\) V0 & V0 \(>\) V3 \\
0.30 & 28.7\% & V3 \(>\) V0 & 47.0\% & V3 \(>\) V0 & V0 \(>\) V3 \\
0.40 & 31.8\% & V3 \(>\) V0 & 50.2\% & V3 \(>\) V0 & V0 \(>\) V3 \\
0.45 & 33.3\% & V3 \(>\) V0 & 52.5\% & V3 \(>\) V0 & V0 \(>\) V3 \\
\bottomrule
\end{tabular}
\begin{tablenotes}\footnotesize
\item \(V0\) is the full SME-only set-aside and \(V3\) is the 10 percent
preference benchmark. The non-pharmaceutical ranking is stable across the
\(\lambda\) grid. In pharmaceuticals, the ranking depends on how the post-policy
SME pool is modeled, not mainly on the marginal cost of public funds.
\end{tablenotes}
\end{threeparttable}
\end{table}

\subsection{Annual scaling}

Annual scaling is kept out of the main text because it is not part of the core
identification exercise. For reference, Group~65 annual reference outlays are
R\$\welfRefOutlayNpYr{} million in non-pharmaceuticals and
R\$\welfRefOutlayPharmaYr{} million in pharmaceuticals. Applying the structural
welfare losses and adherence rates between \welfAdherenceLoPct{} and
\welfAdherenceHiPct{} percent gives an annual welfare-cost range of
R\$\welfRangeLoBRL--\welfRangeHiBRL{} million, or
US\$\welfRangeLoUSD--\welfRangeHiUSD{} million at an exchange rate of
\fxRate{} R\$/US\$. This scaling is useful for magnitude, but it should not be
read as a claim about all procurement categories or about dynamic long-run SME
capacity.

\bibliographystyle{chicago}
\bibliography{References}

\end{document}
